\documentclass[11pt]{article}
\usepackage[margin=1in]{geometry}
\usepackage{amsmath,amssymb,amsthm}
\usepackage{graphicx}
\usepackage[hidelinks]{hyperref}

\newtheorem{theorem}{Theorem}
\newtheorem{definition}{Definition}

\newcommand{\N}{\mathbb N}
\newcommand{\fin}{\mathrm{Fin}}
\newcommand{\restr}{\upharpoonright}
\newcommand{\cl}{\operatorname{cl}}
\newcommand{\inter}{\operatorname{int}}

\title{A Branch-Regular Tree Reduction for Large Almost Disjoint Families on \(\mathbb N^*\)}
\author{Run \texttt{fa\_banach\_001}, lane 13}
\date{2026-06-16}

\begin{document}
\maketitle

\section*{Status}

Conditional reduction, likely valid. The proof below is complete assuming the
explicit tree principle in Definition 1. The packet does not prove that this
tree exists in ZFC.

\section*{Source Question}

Lambie-Hanson--Schrittesser \cite{LHS2026} define an almost disjoint family on
a topological space \(X\) to be a family of noncompact regular open subsets of
\(X\) whose pairwise intersections are compact. Question 4.3 asks whether ZFC
proves the existence of such a family on \(\mathbb N^*\) of cardinality
\(2^{\aleph_1}\), or even \(2^{2^{\aleph_0}}\).

\begin{center}
\includegraphics[width=\textwidth]{figures/open_problem_crop.png}
\end{center}

\section*{Conditional Dependencies}

The only unproved dependency is the following tree principle.

\begin{definition}[Binary branch-regular Hausdorff tree]
A binary branch-regular Hausdorff tree in \(\mathcal P(\omega)/\fin\) is a
family
\[
  \langle a_s:s\in 2^{<\omega_1}\rangle
\]
of infinite subsets of \(\omega\) satisfying:
\begin{enumerate}
\item If \(s\subsetneq t\), then \(a_s\subseteq^* a_t\) and
  \(a_t\setminus a_s\) is infinite.
\item If \(s,t\in2^{<\omega_1}\) are distinct, then
  \[
    a_s\cap a_t={}^*a_{s\wedge t},
  \]
  where \(s\wedge t\) is their longest common initial segment.
\item For every \(s\in2^{<\omega_1}\) and \(i\in\{0,1\}\),
  \(a_{s^\frown i}\setminus a_s\) is infinite.
\item For every branch \(f\in2^{\omega_1}\) and every infinite
  \(y\subseteq\omega\), either
  \[
    y\subseteq^* a_{f\restr\alpha}
  \]
  for some \(\alpha<\omega_1\), or there is an infinite \(z\subseteq y\)
  such that
  \[
    z\cap a_{f\restr\alpha}={}^*\emptyset
    \qquad\text{for all }\alpha<\omega_1.
  \]
\end{enumerate}
\end{definition}

The last clause is the branch-regularity clause. It is exactly the place where
the source paper's \(\mathfrak b=2^{\aleph_0}\) hypothesis performs work in
its construction of a stronger \(\omega\)-branching tree.

\section*{Idea of the Proof}

Given a branch \(f\), take the open subset of \(\mathbb N^*\) generated by the
clopens \(N_{a_{f\restr\alpha}}\). Branch-regularity makes this open set
regular open. Strict growth along the branch makes it noncompact. If two
branches split at a node \(r\), the controlled-intersection clause forces the
intersection of the two corresponding open sets to be precisely the clopen
\(N_{a_r}\). Thus the branches of the tree become the desired generalized
almost disjoint family.

\begin{theorem}
If there exists a binary branch-regular Hausdorff tree in
\(\mathcal P(\omega)/\fin\), then there is an almost disjoint family on
\(\mathbb N^*\) of cardinality \(2^{\aleph_1}\).
\end{theorem}

\begin{proof}
For \(a\in[\omega]^\omega\), let
\[
  N_a=\{p\in\mathbb N^*:a\in p\}
\]
be the corresponding basic clopen subset of \(\mathbb N^*\). Let
\(\langle a_s:s\in2^{<\omega_1}\rangle\) be a binary branch-regular Hausdorff
tree. For each branch \(f\in2^{\omega_1}\), define
\[
  U_f=\bigcup_{\alpha<\omega_1}N_{a_{f\restr\alpha}}.
\]
This set is open.

We first show that \(U_f\) is regular open. Let \(p\in\mathbb N^*\setminus
U_f\), and let \(N_y\) be a basic clopen neighborhood of \(p\), with
\(y\in[\omega]^\omega\). Since \(p\notin U_f\), we cannot have
\(y\subseteq^*a_{f\restr\alpha}\) for any \(\alpha<\omega_1\); otherwise
\(a_{f\restr\alpha}\in p\), so \(p\in U_f\). By branch-regularity, choose an
infinite \(z\subseteq y\) such that
\[
  z\cap a_{f\restr\alpha}={}^*\emptyset
  \quad\text{for all }\alpha<\omega_1.
\]
Then \(N_z\subseteq N_y\) and \(N_z\cap U_f=\emptyset\). Thus every basic
neighborhood of every point outside \(U_f\) contains a smaller basic clopen
set disjoint from \(U_f\). Equivalently, no point of
\(\mathbb N^*\setminus U_f\) lies in \(\inter\cl(U_f)\). Since \(U_f\) is
open, it follows that \(U_f=\inter\cl(U_f)\), so \(U_f\) is regular open.

Next \(U_f\) is not compact. The family
\[
  \{N_{a_{f\restr\alpha}}:\alpha<\omega_1\}
\]
is an increasing open cover of \(U_f\). No finite subfamily covers \(U_f\):
if \(\alpha<\omega_1\), then
\(a_{f\restr(\alpha+1)}\setminus a_{f\restr\alpha}\) is infinite, so there is
a nonprincipal ultrafilter containing this difference; it lies in
\(N_{a_{f\restr(\alpha+1)}}\setminus N_{a_{f\restr\alpha}}\). Hence \(U_f\)
is noncompact.

Now let \(f\neq g\) be two branches. Let \(\delta<\omega_1\) be the first
ordinal with \(f(\delta)\neq g(\delta)\), and put \(r=f\restr\delta
=g\restr\delta\). We claim that
\[
  U_f\cap U_g=N_{a_r}.
\]
The inclusion \(N_{a_r}\subseteq U_f\cap U_g\) follows from
\(a_r\subseteq^*a_{f\restr\alpha}\) and
\(a_r\subseteq^*a_{g\restr\beta}\) for all sufficiently large
\(\alpha,\beta>\delta\).

Conversely, if \(p\in U_f\cap U_g\), choose \(\alpha,\beta<\omega_1\) such
that \(a_{f\restr\alpha}\in p\) and \(a_{g\restr\beta}\in p\). Then
\[
  a_{f\restr\alpha}\cap a_{g\restr\beta}\in p.
\]
By the controlled-intersection clause, this intersection is equal modulo
finite to \(a_{(f\restr\alpha)\wedge(g\restr\beta)}\), and the latter is
equal modulo finite to \(a_r\) after increasing \(\alpha,\beta\) past
\(\delta+1\). Since \(p\) is nonprincipal, \(a_r\in p\), so \(p\in N_{a_r}\).
Thus \(U_f\cap U_g=N_{a_r}\), a clopen and therefore compact subset of
\(\mathbb N^*\).

Finally, the sets \(U_f\) are distinct for distinct branches. If \(f,g\) split
at \(r\) and \(s=r^\frown f(\delta)\), then
\(a_s\setminus a_r\) is infinite and is almost disjoint from every node on the
branch \(g\) above \(r\). A nonprincipal ultrafilter containing
\(a_s\setminus a_r\) belongs to \(U_f\setminus U_g\). Hence
\(\{U_f:f\in2^{\omega_1}\}\) has cardinality \(2^{\aleph_1}\).

We have produced \(2^{\aleph_1}\) many noncompact regular open subsets of
\(\mathbb N^*\) with compact pairwise intersections, as required.
\end{proof}

\section*{Verification and Remaining Work}

The proof above uses the tree dependency only in three places: branch
regularity gives regular openness, strict branch growth gives noncompactness,
and controlled intersections give compact pairwise intersections. Without the
dependency, no conclusion is claimed.

To turn this conditional packet into a full solution of the first part of
Question 4.3, it remains to construct the binary branch-regular Hausdorff tree
in ZFC. The source paper constructs a related stronger tree under
\(\mathfrak b=2^{\aleph_0}\); removing that hypothesis is the open core.

\section*{Literature and Search Notes}

The local run indexes were checked for arXiv:2606.02040 and the existing
continuum-sized partial packet. Web searches were run for Hausdorff gaps,
regular open subsets of \(\omega^*\), almost disjoint families on
\(\omega_1\), and branch/tree formulations in \(\mathcal P(\omega)/\fin\).
The searches did not reveal a ZFC construction of the dependency above. They
did reinforce that Hausdorff and Luzin gaps are the natural nearby tools.

\section*{Human Review Recommendation}

Reviewers should focus on the conditional dependency. The implication from
the tree to the topological family is straightforward; the real mathematical
question is whether the branch-regular tree can be built in ZFC, or whether
its existence implies a cardinal-invariant hypothesis such as a domination
principle.

\begin{thebibliography}{9}

\bibitem{LHS2026}
C. Lambie-Hanson and D. Schrittesser,
\emph{Generalized almost disjoint families and injective Banach spaces},
arXiv:2606.02040, 2026.
\url{https://arxiv.org/abs/2606.02040}

\end{thebibliography}

\end{document}

